% Verbatim copy of /Users/alingalatan/workspaces/grothe/handoff/manuscript/legacy/ch25-parts/85-weighted-full-energy.tex (file dated 2026-10-25). Only this line was added.
\subsection{Visible support and a weighted full-energy response region}
\label{sec:weighted-full-energy-response}
\paragraph{Status and scope (22 October).}
The inequalities below are \st{P,C}: analytic proofs with fresh exact
integer-interval replays of eight scalar integrals, including their infinite
tails. They concern actual bounded responses in every finite dimension,
with all sectors, arbitrary dependence and phases. No cubic moment bound,
raw K, buffer, or tensor-orientation hypothesis is used. This is an
additional region for the fixed witness $\epsilon=1891/2500$, not a
complete row cover or a new global endpoint. The prior endpoint integration
is a separate task.

Use the notation of \S\ref{sec:anisotropic-response-cap}: $Ph=tZ$,
$F=\epsilon Qk$, $p=1/\epsilon$, $T=|Z|^2$, and
\[
 F=q+(T-1)r_A+Z^2r_B/\sqrt2+ZY+\bar ZW+\nu H(Z).
\]
In particular $e=\|q\|^2$ is remote energy, whereas $\|F\|^2$ includes
every sector. The objective is
$\mathcal N=t^2+p\|F\|^2-\|Pk\|^2-\epsilon\|Qh\|^2$.

\subsubsection{A closed-form visible support constraint [P,C]}
For $r\ge0$ define
\[
 I(r)=\E\sqrt T\,|T-r|
 =(3/2-r)\frac{\sqrt\pi}{2}
       +2\int_0^r\sqrt s(r-s)e^{-s}\,ds.
\]
Since $|h|\le1$ and $\E[\bar Zh]=t$,
\[
 m:=\Re\E[\bar Z(T-2)h]
 \le I(r)+(r-2)t.
\]
The visible cubic $H(Z)=Z(T-2)/\sqrt2$ has unit norm. Thus
\begin{equation}\label{eq:visible-support-penalty}
 \boxed{\quad
 \|Pk\|^2+\epsilon\|Qh\|^2
 \ge\frac\epsilon2\bigl((2-r)t-I(r)\bigr)_+^2.
 \quad}
\end{equation}
This uses the visible average before squaring. It is a necessary condition
on actual responses, with no restriction on $F$.

For $r=1/10$ an alternating series gives
\[
 I(r)=(3/2-r)\frac{\sqrt\pi}{2}
 +2r^{5/2}\sum_{n\ge0}
 \frac{(-r)^n}{n!(n+3/2)(n+5/2)}
 <\frac{124233393}{10^8}.
\]
Machin's formula for $\pi$, integer square roots and forty alternating
terms certify this upper bound. At $t=1057/1250$,
\eqref{eq:visible-support-penalty} forces a penalty greater than
$0.05019429$. Consequently $\|F\|^2\le0.03662$ already implies
$\mathcal N<0.713260$. The theorem below improves this cutoff and
retains the sector weights.

This constraint is not redundant with the degree-two multiplication
matrices in \S\ref{sec:response-moment-compressions}. For example take
formal moments on two coordinates
\[
 p^\pm=(t,0),\qquad m^\pm_{00;0}=-5/16,\qquad
 m^\pm_{11;1}=\pm1/3,
\]
and set other cubic moments to zero. Both displayed multiplication
contractions and truncated Parseval pass: their squared Gram diagonals are
\[
 \bigl(t^2+(t-5/16)^2,\ t^2+1/9\bigr),\qquad
 \bigl((2t-5/16)^2/2,\ t^2+1/18\bigr),
\]
all below $1$ at $t=1057/1250$. Their formal objective is
$t^2+\epsilon/18-25\epsilon/512>500/701$.
But their visible moment $-5/16$ violates the support bound above.
These are fake moment data, not actual responses or a counterexample to K.
The exact checker regenerates both matrices and verifies these comparisons.

\subsubsection{The weighted full-energy theorem [P,C]}
Define
\begin{equation}\label{eq:weighted-full-energy-definition}
\begin{aligned}
 \mathcal V_*={}&e+\frac7{10}(\|r_A\|^2+\|r_B\|^2)
  +\frac34\|Y_R\|^2+\frac9{10}\|Y_I\|^2+\frac56\|W\|^2\\
 &+\frac7{20}(\Re\nu)^2+\frac{61}{100}(\Im\nu)^2.
\end{aligned}
\end{equation}
Every weight lies in $[0,1]$, so $\mathcal V_*\le\|F\|^2$.
It is neither the remote variance $e$ nor the earlier adjusted energy
$\mathcal E_{\rm adj}$.

\begin{theorem}[\st{P,C}: full-energy response cap]
\label{thm:weighted-full-energy-response}
Every actual response at $\epsilon=1891/2500$ satisfies
\begin{equation}\label{eq:weighted-full-energy-cap}
 \boxed{\quad
 \mathcal N\le t^2-\frac{1401}{2500}t
              +\frac{1701143}{4000000}+\frac54\mathcal V_*.
 \quad}
\end{equation}
Consequently:
\begin{align}
 0\le t\le851/1000,\quad\mathcal V_*\le13/400
 &\Longrightarrow\quad\mathcal N\le0.71321135,
 \label{eq:weighted-full-energy-strip}\\
 t=1057/1250,\quad\mathcal V_*\le234/6250
 &\Longrightarrow\quad\mathcal N\le0.71325087.
 \label{eq:weighted-full-energy-peak}
\end{align}
Both caps are strictly below $B_*=35663/50000<500/701$.
On $0.84\le t\le0.851$, the tapered condition
\begin{equation}\label{eq:weighted-full-energy-taper}
 \mathcal V_*\le\frac{3253}{100000}
                +\frac{113}{125}\left(\frac{851}{1000}-t\right)
\end{equation}
gives $\mathcal N\le0.71324885<B_*$.
\end{theorem}

\begin{proof}
Use the general weighted two-disk bound \eqref{eq:aniso-general}, with
\[
 a=\frac3{2000},\quad b=\frac{2787}{20000},\quad
 s_0=\frac{47}{500000},\quad d=0,\qquad
 \lambda(T)=\sqrt{T(a+bT)^2+s_0}.
\]
These are free dual parameters, not a change to the witness or the
Conjecture-K profile. Here $2(a+2b)=1401/2500$.
The same eight integrals $J_j,V_j$ and real Gram matrix $K$ as in
\eqref{eq:aniso-gram} give the following directed bounds:
\[
\begin{array}{c|c}
 \text{coefficient of the sector energy in the upper bound}
      &\text{certified upper bound}\\\hline
 p-p^2/C_0&5/4\\
 pI-p^2K^{-1}&(7/8)I\\
 p-p^2/V_1&15/16\\
 p-p^2/J_1&9/8\\
 p-p^2/C_1&25/24\\
 p-p^2/K_H(V)&7/16\\
 p-p^2/K_H(J)&61/80
\end{array}
\qquad
 C_\lambda<\frac{1701143}{4000000}.
\]
All inverse quantities are certified positive. For the A/B comparison,
the difference of the two real $2\times2$ matrices has positive diagonal
and determinant; no phase is suppressed. Multiplying the weights in
\eqref{eq:weighted-full-energy-definition} by $5/4$ gives precisely
the seven rational bounds. Substitution proves
\eqref{eq:weighted-full-energy-cap}.

Unlike the earlier certificate, the sector coefficients here need not be
negative. Their joint weighted energy is retained instead of bounding
each sector independently without an energy budget. The $\|Pk\|^2$
penalty is already used by the Gram correction and cannot be used again.
The compatible square
$\epsilon\|Qh+pbZ(T-2)\|^2$ may still be subtracted from the displayed
upper bound. No such square is needed for the numerical consequences.

The integration proof is the exact one in
\S\ref{sec:anisotropic-response-cap}, with these new parameters:
integer intervals at scale $10^{40}$, fourth-derivative Simpson bounds
on $[0,40]$, and explicit exponential tails. No stored integral or
floating-point decision enters the replay. The default replay uses
2,723 panels; the finer replay uses 7,770 panels and a hundredfold smaller
error target. The finer replay is not a separate analytic assumption.
Finally the right side is a convex quadratic in $t$, even after inserting
the linear taper. Evaluation at the interval endpoints gives the stated
caps by exact rational arithmetic.
\end{proof}

\paragraph{The new necessary condition on a survivor.}
For the design cap $B_*$, an actual response with $\mathcal N>B_*$ must
satisfy
\begin{equation}\label{eq:weighted-full-energy-survivor}
 \boxed{\quad
 \mathcal V_*>\frac45\left\{B_*-t^2+\frac{1401}{2500}t
                          -\frac{1701143}{4000000}\right\}.
 \quad}
\end{equation}
At the peak this threshold is exactly $0.037447304$, and at $t=0.851$
it is $0.03253892$. These conditions must be intersected with the earlier
remote-energy, adjusted-energy, support and polynomial-moment constraints.
They have not been proved jointly inconsistent. In particular the full
energy permitted by the older feasibility bounds is larger than this
new cutoff. A response below one of these caps is covered; a response
above the energy cutoff is not thereby feasible or a violation.

\paragraph{Replay and porting.}
The self-contained proof and standard-library Python checkers are in
\pathref{/Users/alingalatan/Documents/Codex/2026-10-10/plea/outputs/response-attack-20261022/}.
Run \code{python3 -I verify\_all.py}, optionally with \code{--fine}, from
any working directory. The numerical theorem is ready to port using the
existing two-disk and weighted-Gram lemmas; it has not been formalised
in Lean in this round. Float optimisation files in that directory are
diagnostics only and are not inputs to the checker.

\subsubsection{An exact limit of the proportional-support ansatz [P,C]}
\label{sec:proportional-support-obstruction-20261022}
The full joint-support inequality can be optimised analytically in its
scalar multiplier. Put $R=\sqrt\pi/2$,
$\theta=(2R+(3\pi/32)(\eta^{-1}-1))^{-1}$ and $b_0=(1-\eta)/4$.
With $x=P_Xk$, $y=pr_\parallel$, its constraint is
\[
 \theta\|x-a_1\|^2+b_0\|y-a_2\|^2
 \le d:=D-\theta m^2-p^2E_0/4.
\]
For $\mu>1/(b_0p)$, simultaneous square completion bounds the remaining
cost $-\|x\|^2+\|y\|^2/p+pE_0$ by
\[
 \mathcal F(\mu)=\mu d-\frac{\mu\theta\|a_1\|^2}{1+\mu\theta}
 +\frac{\mu b_0\|a_2\|^2}{\mu b_0p-1}+pE_0.
\]
Its second derivative is nonnegative, so this scalar optimisation needs
only a root or boundary comparison, not a nonconvex search.

In the diffuse Gaussian-grade model $q=\sqrt e\,W$, assign $W$ grade
three, independently of the first-grade variables. For proportional
tests $v=\lambda q$, all complex $\lambda$ are allowed and the
first-grade centres vanish. Writing $z=\sqrt e\lambda$ gives
\[
 D(z)=R\sqrt{1+|z|^2}-p\sqrt e\Re z-t,\qquad
 m(z)^2=|\nabla D(z)|^2.
\]
The Hessian of $D$ is bounded above by $RI$. Therefore
$D-D_{\min}\ge m^2/(2R)$, where
$D_{\min}=\sqrt{R^2-p^2e}-t$. When $D_{\min}\ge0$ and $E_0=0$,
\[
 \inf_{\lambda,\eta,\mu}\mathcal F(\mu)=4\epsilon D_{\min}.
\]
The same relaxed vector $x=0$, $\|y\|^2=4D_{\min}$ satisfies
\emph{every} proportional-test cut simultaneously. Thus combinations
of those cuts do not repair this relaxation either.
At $t=1057/1250$, $e=1/25$, the cost is strictly between
$0.0007977136526283$ and $0.0007977136526285$, whereas
$B_*-B_0(t)=10895563/94550000000<0.000115236$.
This is a limitation of the proportional-test relaxation, not an actual
bounded response, a finite-dimensional cubic counterexample, or a failure
of K. The Gaussian-grade family itself is already covered by other
theorems. The full proof and exact replay, including the scalar
square-completion identities, are in the package's \code{support/}
directory. Nonproportional tests and simultaneous response constraints
remain available.

\paragraph{IDE replay and review (22 October).} \code{verify\_all.py --fine} PASSES from an unrelated directory under \code{python3 -I} in 34\,s (7\,770 accepted panels;
package and log in \pathref{gpt-reviews/weighted-energy-20261022/}). The IDE re-derived by hand: the two-disk identity (the gap is $(c_0-xa_0)^2+x^2c_0^2(1-a_0^2)$),
its extension to disks by convexity, the split $L_0=\ell_0Z+B$ with $\ell_0=a+2b$ and $2pb^2=p\|B\|^2$, and the Gram entries of every sector under $M^{-1}u=c\,u+d\,e^{2i\arg Z}\bar u$:
$C_0$ for the remote cubic (no first-chaos correction, since $PM^{-1}q=0$), the $2\times2$ block with partner $Z^2\bar X/\sqrt2$ and the $1/(1+C_0)$ correction,
$V_1$/$J_1$ for $\Re Y$/$\Im Y$, $C_1$ for $W$, and $K_H$ with its $1/(1+V_1)$ visible correction. The sector weights are the certified coefficients divided by $5/4$.
The visible support lemma was checked numerically ($I(1/10)=1.24234$, $m\le-0.3643$, $J>0.0502$). Not checked line by line: the finite-rank inverse formula for $M+P$ across all sectors simultaneously
(each sector was checked separately; the sectors are mutually orthogonal under $M^{-1}$ by charge and remote-degree counting).
\emph{Orientation from the IDE's asymptotic W-model map (attack step 1, \pathref{attack-20261022/})}: at $t=0.8456$ the largest payoff of actual responses $\ph(Za(T)\pm W)$
with remote variance $e\le0.035$ is $0.7091$, and with $e=0.041$ at $t=0.844$ it is $0.71314$. So, in that model, the intermediate-variance band is at least $4\cdot10^{-3}$ below $500/701$; the
loss is in the certificates, not in the responses.
